\chapter{Next.js App Router, Routing \& the Proxy Layer}
\label{ch:proxy}

\section{App Router conventions used in this project}

Every folder under \pathname{src/app} is a URL segment; a folder becomes a real page only when it contains a \pathname{page.tsx}. Three conventions are used throughout OLK:

\begin{itemize}
  \item \textbf{Route groups} — \pathname{(dashboard)} wraps most of the app in a shared \pathname{layout.tsx} without adding \code{/dashboard} to any URL. \pathname{/browse} is served by \pathname{src/app/(dashboard)/browse/page.tsx}.
  \item \textbf{Dynamic segments} — \pathname{[id]} folders (\pathname{messages/[id]}, \pathname{clubs/[id]}, \pathname{user/[id]}) receive the segment value as a route param.
  \item \textbf{Nested layouts} — \pathname{admin/layout.tsx} wraps every \pathname{admin/*} page a second time, inside the outer \pathname{(dashboard)/layout.tsx}, so an admin page is rendered inside two layouts stacked: dashboard shell, then admin nav.
\end{itemize}

\section{The Proxy layer — the rename, applied}

\begin{olknote}
\textbf{Next.js 16 renamed Middleware to Proxy.} From the framework's own documentation (\pathname{node\_modules/next/dist/docs/01-app/03-api-reference/03-file-conventions/proxy.md}):

\begin{quote}
\emph{"The \texttt{middleware} file convention is deprecated and has been renamed to \texttt{proxy}."}
\end{quote}

The convention is a \pathname{proxy.ts} file exporting a function named \code{proxy} (or a default export), instead of \pathname{middleware.ts} exporting \code{middleware}. OLK now uses \pathname{src/proxy.ts} — the body of the function is otherwise unchanged from the old \pathname{middleware.ts}, per the official migration guidance (\code{npx @next/codemod@canary middleware-to-proxy .}): only the filename and the exported function's name moved.
\end{olknote}

\section{What the Proxy actually does, line by line}

\begin{lstlisting}[caption={\pathname{src/proxy.ts} — full file}]
import { createServerClient } from '@supabase/ssr'
import type { SupabaseClient } from '@supabase/supabase-js'
import { NextResponse, type NextRequest } from 'next/server'
import { FEATURE_FLAG_KEYS, parseFeatureFlags, type FeatureFlags } from '@/lib/platform-settings'

const PROTECTED_ROUTES = ['/my-books', '/messages', '/requests', '/profile', '/user', '/saved', '/wishlist', '/clubs/create', '/admin']

let cachedFlags: { flags: FeatureFlags; expiresAt: number } | null = null
const FLAGS_TTL_MS = 30_000

async function getFeatureFlags(supabase: SupabaseClient): Promise<FeatureFlags> {
  const now = Date.now()
  if (cachedFlags && cachedFlags.expiresAt > now) return cachedFlags.flags
  const { data } = await supabase.from('platform_settings').select('key, value').in('key', FEATURE_FLAG_KEYS)
  const flags = parseFeatureFlags(data)
  cachedFlags = { flags, expiresAt: now + FLAGS_TTL_MS }
  return flags
}

export async function proxy(request: NextRequest) {
  let supabaseResponse = NextResponse.next({ request })

  const supabase = createServerClient(
    process.env.NEXT_PUBLIC_SUPABASE_URL!,
    process.env.NEXT_PUBLIC_SUPABASE_ANON_KEY!,
    {
      cookies: {
        getAll() { return request.cookies.getAll() },
        setAll(cookiesToSet) {
          cookiesToSet.forEach(({ name, value }) => request.cookies.set(name, value))
          supabaseResponse = NextResponse.next({ request })
          cookiesToSet.forEach(({ name, value, options }) =>
            supabaseResponse.cookies.set(name, value, options)
          )
        },
      },
    }
  )

  const { data: { user } } = await supabase.auth.getUser()
  const { pathname } = request.nextUrl
  const isProtected = PROTECTED_ROUTES.some(route => pathname.startsWith(route))
  const flags = await getFeatureFlags(supabase)

  let profile: { is_banned: boolean; ban_expires_at: string | null; is_admin: boolean; admin_role: string | null } | null = null
  if (user && (isProtected || flags.maintenance_mode)) {
    const { data } = await supabase
      .from('profiles')
      .select('is_banned, ban_expires_at, is_admin, admin_role')
      .eq('id', user.id)
      .single()
    profile = data
  }

  if (flags.maintenance_mode && pathname !== '/maintenance' && pathname !== '/login' && !profile?.is_admin) {
    const maintenanceUrl = request.nextUrl.clone()
    maintenanceUrl.pathname = '/maintenance'
    return NextResponse.redirect(maintenanceUrl)
  }

  if (!flags.feature_clubs && pathname.startsWith('/clubs')) {
    const browseUrl = request.nextUrl.clone(); browseUrl.pathname = '/browse'; return NextResponse.redirect(browseUrl)
  }
  if (!flags.feature_wishlists && pathname.startsWith('/wishlist')) {
    const browseUrl = request.nextUrl.clone(); browseUrl.pathname = '/browse'; return NextResponse.redirect(browseUrl)
  }
  if (!flags.feature_messages && pathname.startsWith('/messages')) {
    const browseUrl = request.nextUrl.clone(); browseUrl.pathname = '/browse'; return NextResponse.redirect(browseUrl)
  }

  if (!user && isProtected) {
    const loginUrl = request.nextUrl.clone()
    loginUrl.pathname = '/login'
    return NextResponse.redirect(loginUrl)
  }

  if (user && isProtected) {
    if (profile?.is_banned) {
      const expired = profile.ban_expires_at && new Date(profile.ban_expires_at) < new Date()
      if (!expired && !pathname.startsWith('/profile')) {
        const bannedUrl = request.nextUrl.clone()
        bannedUrl.pathname = '/profile'
        return NextResponse.redirect(bannedUrl)
      }
    }

    if (pathname.startsWith('/admin') && !profile?.is_admin) {
      const browseUrl = request.nextUrl.clone()
      browseUrl.pathname = '/browse'
      return NextResponse.redirect(browseUrl)
    }
  }

  return supabaseResponse
}

export const config = {
  matcher: ['/((?!_next/static|_next/image|favicon.ico|.*\\.(?:svg|png|jpg|jpeg|gif|webp)$).*)'],
}
\end{lstlisting}

Five responsibilities, in the order they execute:

\begin{enumerate}
  \item \textbf{Session refresh} — the \code{cookies.getAll}/\code{setAll} dance exists purely so \code{@supabase/ssr} can rotate the auth token cookies on every request. This runs unconditionally, for \emph{every} request matched by \code{config.matcher} (i.e. everything except static assets), whether or not the route is protected.
  \item \textbf{Feature flags} — \code{platform\_settings} (Chapter~\ref{ch:admin}) is read once per (cached, 30-second-TTL) window via \pathname{src/lib/platform-settings.ts}, shared with the \pathname{(dashboard)/layout.tsx} server component that hides nav links for disabled features.
  \item \textbf{Maintenance mode} — if \code{maintenance\_mode} is on, every non-admin visitor (whether logged in or not) is redirected to \pathname{/maintenance}, except \pathname{/login} itself (so an admin can still sign in) and \pathname{/maintenance}.
  \item \textbf{Per-feature route gates} — \code{feature\_clubs}/\code{feature\_wishlists}/\code{feature\_messages} being off redirects their entire route tree (\pathname{/clubs}, \pathname{/wishlist}, \pathname{/messages}) to \pathname{/browse}, independent of whether the route is otherwise "protected."
  \item \textbf{Unauthenticated, banned-user, and admin gates} — unchanged from before: redirect a session-less visitor away from \code{PROTECTED\_ROUTES} to \pathname{/login}; for a logged-in user on a protected route, bounce a currently-banned user to \pathname{/profile}, and bounce any non-admin away from \pathname{/admin}.
\end{enumerate}

\begin{olknote}
Look again at \code{PROTECTED\_ROUTES}: it does \emph{not} include \pathname{/browse} or \pathname{/notifications} or \pathname{/clubs} (only \pathname{/clubs/create}). This is deliberate — those pages are meant to be at least partially visible to guests — but it means any new page added under \pathname{(dashboard)} that assumes a logged-in \code{user} object without checking for \code{null} itself will crash for anonymous visitors unless it is also added to this list, or unless the page guards \code{auth.getUser()} being null on its own. There is no automatic enforcement tying new routes to this list — it is a manually maintained array. The feature-flag route gates above are a separate, independent set of checks — a route can be flag-gated without being in \code{PROTECTED\_ROUTES}, and vice versa.
\end{olknote}

\begin{olkhowto}
Adding a new page that should require login (say \pathname{/my-loans}) means touching three places, none of which reference each other:
\begin{enumerate}
  \item \pathname{src/app/(dashboard)/my-loans/page.tsx} — the page itself, inside the \pathname{(dashboard)} route group so it inherits the shared shell.
  \item \code{PROTECTED\_ROUTES} in \pathname{src/proxy.ts} — add \code{'/my-loans'}, or an anonymous visitor reaches the page and your code has to handle a null \code{user} entirely on its own.
  \item \code{NAV\_ITEMS} in \pathname{src/components/dashboard-sidebar.tsx} — add \code{\{ href: '/my-loans', label: '...' \}}, or the page exists but nothing in the app links to it.
\end{enumerate}
None of these three will error or warn if you forget one — the page will just be unprotected, or unprotected-but-unlisted, or listed-but-crashing for guests. The habit from Chapter~\ref{ch:wishlists-clubs} applies here too: visit your new URL directly, logged out, before considering the page done.
\end{olkhowto}

\section{The \code{matcher} config}

The exported \code{config.matcher} is a single regex excluding \pathname{\_next/static}, \pathname{\_next/image}, \pathname{favicon.ico}, and common image extensions — meaning the Proxy runs on literally every other request, including API routes and RSC payload fetches, not just page navigations. This is why the session-refresh logic at the top has to be cheap and unconditional: it runs far more often than the redirect logic below it.
